\begin{textsample}{NEG Dim 1 – pa – Score 7.00 – pa/pa\_em\_73.txt}  \label{ex:f1_neg_259}
Quadro 22 : Estrutura das ementas dos Campos de Saberes e Práticas Eletivos Elementos 1.
Título do campo de saberes : Este deve ser atrativo ( linguagem das juventudes ). 2.
Princípio(s) curricular(es) norteador(es) : - Respeito às diversas culturas amazônicas e suas inter-relações no espaço e no tempo ; - Educação para a sustentabilidade ambiental, social e econômica ; - Interdisciplinaridade no processo ensino-aprendizagem.
Os princípios acima listados são fundamentos que norteiam ( orientam ) o currículo na perspectiva da formação humana integral, trazendo para as discussões o desenvolvimento dos processos criativos e produtivos, advindos dos distintos grupos sociais que compõem a Amazônia, considerando a diversidade ambiental, social, econômica e cultural. 3.
Eixos estruturantes : - Investigação científica ; - Processos criativos ; - Mediação e intervenção sociocultural ; - Empreendedorismo social.
Estes eixos conectam as experiências pedagógicas com a realidade cotidiana dos estudantes, favorecendo o desenvolvimento de habilidades importantes para a formação humana integral. 4.
Carga horária : 40h semestrais. 5.
Área(s) de conhecimento(s) : Aponta as áreas Linguagens e suas tecnologias, Matemática e suas tecnologias, Ciências da Natureza e suas tecnologias, Ciências Humanas Sociais Aplicada.
Caso o campo eletivo seja vinculado à EPT, deve -se informar que se trata de um campo cuja itinerância integra a área e a EPT. 6.
Competência específica de área : Referem -se às competências das áreas de conhecimento apresentadas na nucleação da formação geral básica.
Deve -se indicar o seu código alfanumérico.
Ex. : EM13MAT201 ( segunda competência da área de matemática ). 7.
Habilidades a serem trabalhadas : Refere -se às habilidades associadas às competências gerais da BNCC e às habilidades específicas das itinerâncias, conforme determina a Portaria MEC nº 1.432/2018.
Deve -se indicar o seu código alfanumérico.
Ex. : ( EMIFLGG01 ) – 1ª habilidade específica da itinerância na área de Linguagens e suas Tecnologias. 8.
Objetos de conhecimento a serem aprofundados : Apresenta a lista de objetos de conhecimento possíveis para serem aprofundados nos campos eletivos e associados às habilidades. 9.
Bibliografia básica : Lista de materiais bibliográficos consultados pelo professor para subsidiar a prática docente no campo eletivo.
Fonte : Elaboração com base nos estudos e sistematizações do ProBNCC – etapa ensino médio ( 2019 ).
Assim, os campos de saberes e práticas de ensino eletivos são unidades curriculares flexíveis que possibilitam a integração curricular entre os objetos de ensino que compõem uma área de conhecimento e a EPT.
Dessa forma, o DCEPA – etapa ensino médio apresenta em anexo um conjunto de ementas para subsidiar a elaboração de campos eletivos para serem disponibilizados nas escolas e suas respectivas ementas, que se correlacionarão com os princípios curriculares norteadores, os eixos estruturantes, as competências específicas de área, as habilidades específicas das itinerâncias e os objetos de conhecimento.
Portanto, os campos eletivos são oportunidades de aprendizagens diversas que aliam os interesses dos estudantes, o aprofundamento dos objetos de conhecimento de uma área e seus projetos de vida, que serão discutidos na continuidade desta seção.

% matched lemmas: 
\end{textsample}
